\pagestyle{empty}
\tight
\begin{tblr}{width=\textwidth,colspec={|Q[c,m,1.9cm]|X[l,m]|},hlines,vlines,hspan=minimal,row{1}={bg=headblue},rowsep=1pt,colsep=3pt}
\SetCell{c,m}\hfil\logo\hfil & \textbf{POLITEKNIK NEGERI MALANG}\par\textbf{JURUSAN TEKNOLOGI INFORMASI}\par\textbf{PROGRAM STUDI: D4 TEKNIK INFORMATIKA} \\
\SetCell[c=2]{c} \textbf{RENCANA PEMBELAJARAN SEMESTER (RPS)} & \\
\end{tblr}
\begin{longtblr}{width=\textwidth,colspec={|Q[l,m,1.9cm]|X[0.85,l,m]|X[1.45,l,m]|X[0.75,l,m]|X[1.15,l,m]|},hlines,vlines,hspan=minimal,rowsep=1pt,colsep=2pt,row{1,3}={bg=headgray,font=\bfseries}}
MATA KULIAH & KODE & BOBOT (SKS)/jam & SEMESTER & TGL. PENYUSUNAN \\
Basis Data & RTI252006 & 2 SKS/4 jam & 2 & 29 Januari 2026 \\
OTORISASI & Dosen Pengembang RPS & Koordinator RMK & \SetCell[c=2]{l} Ka PRODI &  \\
 & Candra Bella Vista, S.Kom., M.T.\par Prof. Ir. Yan Watequlis Syaifudin, ST., M.MT., Ph.D. & Candra Bella Vista, S.Kom., M.T. & \SetCell[c=2]{l}{\vspace{8mm}\par Dr. Ely Setyo Astuti, S.T., M.T.} &  \\
\SetCell[r=10]{l} Capaian Pembelajaran (CP) & \SetCell[c=4]{l,bg=headgray,font=\bfseries} Capaian Pembelajaran Lulusan Program Studi (CPL-Prodi) &  &  &  \\
 & CPL02 & \SetCell[c=3]{l} Menguasai konsep teoritis bidang pengetahuan mengenai berbagai prinsip rekayasa sistem/perangkat lunak dalam merancang solusi aplikasi multi-platform yang relevan dengan kebutuhan stakeholder. &  &  \\
 & CPL10 & \SetCell[c=3]{l} Memiliki kompetensi untuk menganalisis persoalan kompleks untuk mengidentifikasi solusi bidang informatika/ilmu komputer dengan mempertimbangkan wawasan ilmu multidisiplin. &  &  \\
 & \SetCell[c=4]{l,bg=headgray,font=\bfseries} Capaian Pembelajaran mata kuliah (CPL-MK) &  &  &  \\
 & CPMK0203 & \SetCell[c=3]{l} Mampu merancang skema dan arsitektur manajemen data yang efektif dan terstruktur untuk mendukung solusi aplikasi multi-platform dengan mempertimbangkan integritas data, kebutuhan informasi stakeholder, serta efisiensi dalam operasi query dan skalabilitas sistem. &  &  \\
 & CPMK1001 & \SetCell[c=3]{l} Mampu menganalisis permasalahan pengelolaan data dan informasi yang bersifat kompleks dengan wawasan multidisiplin untuk merancang serta mengimplementasikan solusi berbasis basis data yang sesuai dengan kebutuhan organisasi atau bisnis. &  &  \\
 & \SetCell[c=4]{l,bg=headgray,font=\bfseries} Kemampuan akhir tiap tahapan belajar (Sub-CPMK) &  &  &  \\
 & SCPMK0203-01501 & \SetCell[c=3]{l} Mampu memahami konsep data, basis data, dan basis data relasional, serta merancang model dan skema basis data secara konseptual dan logikal dengan memperhatikan integritas data dan konsep normalisasi. &  &  \\
 & SCPMK1001-01501 & \SetCell[c=3]{l} Mampu mengimplementasikan skema basis data dan menyusun query SQL untuk mendefinisikan, memanipulasi, dan menampilkan data sesuai dengan studi kasus. &  &  \\
 & SCPMK1001-01502 & \SetCell[c=3]{l} Mampu menganalisis kebutuhan data, merancang skema basis data, dan menyusun query SQL secara efisien untuk mengelola basis data sesuai dengan studi kasus. &  &  \\
\SetCell[r=2]{l} Penilaian & \SetCell[c=4]{l} *Nilai CPMK didapat dari agregasi nilai SUB-CPMK yang dibebankan &  &  &  \\
 & \SetCell[c=4]{l}{\tiny\begin{tblr}{width=\linewidth,colspec={|Q[c,m,0.1\linewidth]|Q[c,m,0.16\linewidth]|X[l,m]|X[l,m]|X[l,m]|X[l,m]|X[l,m]|Q[c,m,0.08\linewidth]|},hlines,vlines,hspan=minimal,rowsep=1pt,colsep=0.7pt,row{1,2}={bg=headgray,font=\bfseries}}
Kode CPMK & Kode SCPMK & \SetCell[c=5]{c} Bobot per Bentuk Penilaian & & & & & Total Bobot \\
 & & Tugas & Kuis & UTS & Case Method & UAS & \\
\SetCell[r=1]{c} CPMK0203 & SCPMK0203-01501 & 10\%: Tugas 1--7 konsep dan perancangan basis data & 10\%: Kuis 1 materi mg 1--3 & 30\%: UTS materi mg 1--8 & & & 50\% \\
\SetCell[r=2]{c} CPMK1001 & SCPMK1001-01501 & 10\%: Tugas 8--9 implementasi SQL & 10\%: Kuis 2 materi mg 10--13 & & & & 20\% \\
 & SCPMK1001-01502 & & & & 10\%: studi kasus perancangan-implementasi mg 15--16 & 20\%: UAS SQL studi kasus & 30\% \\
\SetCell[c=7]{r}\textbf{Total Per Penilaian} & & & & & & & \textbf{100\%} \\
\end{tblr}} &  &  &  \\
Diskripsi Singkat Mata Kuliah & \SetCell[c=4]{l} Mata kuliah Basis Data membahas konsep dasar basis data relasional, perancangan basis data menggunakan ERD, pemetaan ke model relasional, serta normalisasi. Mahasiswa juga mempelajari implementasi basis data menggunakan MySQL melalui DDL, DML, dan query SELECT, serta penerapannya dalam studi kasus sesuai kebutuhan organisasi atau bisnis. &  &  &  \\
Bahan Kajian / Materi Pembelajaran & \SetCell[c=4]{l}\begin{enumerate}\item Pengantar Basis Data\item Entity Relationship Diagram\item Contextual Data Model\item Physical Data Model dan Relational Model\item Normalisasi Basis Data\item Pengantar MySQL dan DDL\item DML MySQL\item Query SELECT\item SELECT Multi Table\item Studi Kasus (identifikasi kebutuhan, ERD, pemetaan, dan implementasi basis data)\end{enumerate} &  &  &  \\
\SetCell[r=2]{l} Pustaka & \SetCell[c=4]{l} \textbf{Utama:}\begin{enumerate}\item Puspitasari, D. and Hani'ah, M. 2019. Cara Mudah Merancang Basis Data Relasional. Polinema Pers.\item Fathansyah. 2018. Basis Data Revisi Ketiga. Bandung: Penerbit Informatika.\end{enumerate} &  &  &  \\
 & \SetCell[c=4]{l} \textbf{Pendukung:}\begin{enumerate}\item Silberschatz, A., Korth, H., Sudarshan, S. 2019. Database System Concepts Seventh Edition. McGraw-Hill.\item Dewantara, R., Chirzah, D., Tamsir, N., Rahmayuna, N., Rachman, S., Muhajirin. 2023. Perancangan Basis Data. Klaten: Penerbit Lakeisha.\item Silva, R. 2023. MySQL Crash Course. San Francisco: William Pollock.\end{enumerate} &  &  &  \\
Media Pembelajaran & \SetCell[c=2]{l} \textbf{Software:} MySQL/MariaDB, MySQL Workbench/phpMyAdmin, perangkat pemodelan ERD, LMS. &  & \SetCell[c=2]{l} \textbf{Hardware:} Komputer, laptop, proyektor. &  \\
Nama Dosen Pengampu & \SetCell[c=4]{l}\begin{enumerate}\item Candra Bella Vista, S.Kom., M.T.\item Elok Nur Hamdana, S.T., M.T.\item M. Hasyim Ratsanjani, S.Kom., M.Kom.\item Ridwan Rismanto, S.ST., M.Kom., Ph.D.\item Vit Zuraida, S.Kom., M.Kom.\end{enumerate} &  &  &  \\
Matakuliah Syarat & \SetCell[c=4]{l} - &  &  &  \\
\end{longtblr}
\begin{longtblr}{width=\textwidth,colspec={|Q[c,m,0.06\textwidth]|X[1.35,l,m]|X[1.08,l,m]|X[1.16,l,m]|X[0.7,l,m]|X[1.24,l,m]|X[1.06,l,m]|X[0.98,l,m]|Q[c,m,0.05\textwidth]|},hlines,vlines,hspan=minimal,rowhead=2,row{1,2}={bg=headgreen,font=\bfseries},rowsep=1pt,colsep=1.5pt}
Mgg.\par Ke & Kemampuan Akhir Yang Direncanakan (Sub-CP-MK) & Bahan kajian (Materi Pembelajaran) & Bentuk dan Metode Pembelajaran & Estimasi Waktu & Pengalaman Belajar Mahasiswa & Kriteria \& Bentuk Penilaian & Indikator Penilaian & Bobot Penilaian (\%) \\
(1) & (2) & (3) & (4) & (5) & (6) & (7) & (8) & (9) \\
1 & SCPMK0203-01501 & Pengantar basis data: pengertian, kegunaan, karakteristik, jenis, contoh penerapan; konsep basis data relasional serta prinsip dan tahapan pengembangannya. & Bentuk: ceramah, diskusi, latihan analisis. Metode: Case Based Learning (CBL). Penugasan: Tugas 1 mencari dan mengulas contoh penerapan basis data. & 1 x 2 x 50' tatap muka dan tugas terstruktur; 1 x 2 x 50' tugas mandiri. & Mahasiswa menjelaskan pengertian, kegunaan, karakteristik, dan jenis basis data serta memberi contoh penerapan basis data relasional. & Kriteria: ketepatan dan penguasaan. Bentuk: keaktifan diskusi kelompok; ketepatan jawaban tugas. & Ketepatan mencari contoh penerapan basis data; ketepatan penjelasan komponen basis data. & 1 \\
2 & SCPMK0203-01501 & Pemodelan data: konsep, jenis, dan arsitektur pemodelan data; versi dan komponen ER Diagram; requirement data. & Bentuk: ceramah, diskusi, latihan analisis. Metode: Case Based Learning (CBL). Penugasan: Tugas 2 menganalisis kebutuhan data suatu studi kasus. & 1 x 2 x 50' tatap muka dan tugas terstruktur; 1 x 2 x 50' tugas mandiri. & Mahasiswa menjelaskan konsep, jenis, dan arsitektur pemodelan data serta versi dan komponen ER Diagram, lalu mengidentifikasi kebutuhan data studi kasus. & Kriteria: ketepatan dan penguasaan. Bentuk: keaktifan diskusi kelompok; ketepatan jawaban tugas. & Ketepatan analisis kebutuhan data. & 1.5 \\
3 & SCPMK0203-01501 & Perancangan basis data menggunakan ER Diagram: langkah perancangan, penentuan entity, atribut, relationship, kardinalitas, dan partisipan. & Bentuk: ceramah, diskusi, latihan analisis. Metode: Case Based Learning (CBL). Penugasan: Tugas 3 merancang basis data menggunakan ERD dari hasil analisis kebutuhan data. & 1 x 2 x 50' tatap muka dan tugas terstruktur; 1 x 2 x 50' tugas mandiri. & Mahasiswa menerapkan langkah-langkah perancangan basis data menggunakan ER Diagram dan merancang ERD dari studi kasus yang diberikan. & Kriteria: ketepatan dan penguasaan. Bentuk: keaktifan diskusi kelompok; ketepatan jawaban tugas. & Ketepatan rancangan basis data menggunakan ERD. & 1.5 \\
4 & SCPMK0203-01501 & Kuis 1: materi minggu 1--3. & Ujian teori pilihan ganda dan/atau esai, sifat closed book. & Sesuai jadwal asesmen. & Mahasiswa mengerjakan soal kuis materi minggu 1--3 secara mandiri dan jujur. & Kriteria: ketepatan jawaban. Bentuk: ujian. & Ketepatan jawaban dengan kunci jawaban. & 10 \\
5 & SCPMK0203-01501 & Pemetaan ERD ke model relasional: pemetaan entity, atribut, relationship, kardinalitas, dan partisipan; contoh pemetaan. & Bentuk: ceramah, diskusi, latihan analisis. Metode: Case Based Learning (CBL). Penugasan: Tugas 4 memetakan ERD yang diberikan ke model relasional. & 1 x 2 x 50' tatap muka dan tugas terstruktur; 1 x 2 x 50' tugas mandiri. & Mahasiswa menjelaskan algoritma mapping ER Diagram ke model relasional dan melakukan pemetaan berdasarkan algoritma yang diberikan. & Kriteria: ketepatan dan penguasaan. Bentuk: keaktifan diskusi kelompok; ketepatan jawaban tugas. & Ketepatan pemetaan komponen ERD ke dalam model relasional. & 1.5 \\
6 & SCPMK0203-01501 & Praktik pemetaan ERD hasil rancangan ke model relasional dan penilaian kesesuaian model relasional dengan requirement data. & Bentuk: ceramah, diskusi, latihan analisis. Metode: Case Based Learning (CBL). Penugasan: Tugas 5 memetakan ERD rancangan sendiri ke model relasional. & 1 x 2 x 50' tatap muka dan tugas terstruktur; 1 x 2 x 50' tugas mandiri. & Mahasiswa memetakan ERD hasil rancangan sebelumnya ke model relasional dan menilai kesesuaiannya dengan requirement data. & Kriteria: ketepatan dan penguasaan. Bentuk: keaktifan diskusi kelompok; ketepatan jawaban tugas. & Ketepatan model relasional dan kesesuaiannya dengan requirement data. & 1.5 \\
7 & SCPMK0203-01501 & Normalisasi basis data: pengertian, tujuan, manfaat, tahapan normalisasi, dan pembentukan bentuk normal 1 (1NF). & Bentuk: ceramah, diskusi, latihan analisis. Metode: Case Based Learning (CBL). Penugasan: Tugas 6 melakukan normalisasi basis data pada kasus yang diberikan. & 1 x 2 x 50' tatap muka dan tugas terstruktur; 1 x 2 x 50' tugas mandiri. & Mahasiswa menjelaskan pengertian, tujuan, manfaat, dan tahapan normalisasi serta membentuk 1NF dari tabel dan data yang diberikan. & Kriteria: ketepatan dan penguasaan. Bentuk: keaktifan diskusi kelompok; ketepatan jawaban tugas. & Ketepatan menjelaskan tahapan normalisasi; ketepatan hasil normalisasi. & 1.5 \\
8 & SCPMK0203-01501 & Normalisasi basis data lanjut: pembentukan bentuk normal 2 (2NF) dan 3 (3NF); contoh normalisasi 1NF--3NF. & Bentuk: ceramah, diskusi, latihan analisis. Metode: Case Based Learning (CBL). Penugasan: Tugas 7 merancang model relasional menggunakan metode normalisasi basis data. & 1 x 2 x 50' tatap muka dan tugas terstruktur; 1 x 2 x 50' tugas mandiri. & Mahasiswa melakukan pembentukan bentuk normal 1NF hingga 3NF pada proses normalisasi basis data berdasarkan tabel dan data yang diberikan. & Kriteria: ketepatan dan penguasaan. Bentuk: keaktifan diskusi kelompok; ketepatan jawaban tugas. & Ketepatan hasil normalisasi hingga 3NF. & 1.5 \\
9 & SCPMK0203-01501 & UTS: materi minggu 1--8. & Ujian teori pilihan ganda dan/atau esai, sifat closed book. & Sesuai jadwal asesmen. & Mahasiswa mengerjakan soal UTS materi minggu 1--8 secara mandiri dan jujur. & Kriteria: ketepatan jawaban. Bentuk: ujian. & Ketepatan jawaban dengan kunci jawaban. & 30 \\
10 & SCPMK1001-01501 & Tahapan implementasi basis data: membuat basis data, tabel, atribut, primary key, foreign key; bahasa SQL dan SQL-DDL (CREATE, ALTER, DROP). & Bentuk: ceramah, diskusi, latihan analisis. Metode: Case Based Learning (CBL). Penugasan: Tugas 8 menuliskan perintah DDL untuk mengimplementasikan hasil rancangan basis data. & 1 x 2 x 50' tatap muka dan tugas terstruktur; 1 x 2 x 50' tugas mandiri. & Mahasiswa menjelaskan tahapan implementasi basis data serta menuliskan perintah SQL-DDL untuk mengimplementasikan dan mengelola basis data, tabel, atribut, dan relasi. & Kriteria: ketepatan dan penguasaan. Bentuk: keaktifan diskusi kelompok; ketepatan jawaban tugas. & Ketepatan, kelengkapan, dan kerapian perintah SQL-DDL untuk implementasi rancangan; ketepatan perintah mengubah dan menghapus basis data, tabel, dan relasi. & 2.5 \\
11 & SCPMK1001-01501 & Pengelolaan data pada basis data: penambahan, penghapusan, perubahan data; bahasa SQL-DML (INSERT, DELETE, UPDATE, klausa WHERE). & Bentuk: ceramah, diskusi, latihan analisis. Metode: Case Based Learning (CBL). Penugasan: Tugas 9 menuliskan perintah DML untuk mengelola data yang tersimpan pada basis data. & 1 x 2 x 50' tatap muka dan tugas terstruktur; 1 x 2 x 50' tugas mandiri. & Mahasiswa menjelaskan proses pengelolaan data serta menuliskan perintah SQL-DML untuk menambah, menghapus, dan mengubah data yang tersimpan pada basis data. & Kriteria: ketepatan dan penguasaan. Bentuk: keaktifan diskusi kelompok; ketepatan jawaban tugas. & Ketepatan perintah SQL-DML untuk menambah, menghapus, dan mengubah data sesuai studi kasus. & 2.5 \\
12 & SCPMK1001-01501 & Query data: proses query data; bahasa SQL-DQL (SELECT, klausa WHERE, ORDER BY, GROUP BY, fungsi agregasi SUM, MIN, MAX, AVG). & Bentuk: ceramah, diskusi, latihan analisis. Metode: Case Based Learning (CBL). Penugasan: menuliskan perintah SELECT untuk menampilkan kembali data pada basis data. & 1 x 2 x 50' tatap muka dan tugas terstruktur; 1 x 2 x 50' tugas mandiri. & Mahasiswa menjelaskan proses query data serta menuliskan perintah SQL-DQL untuk menampilkan kembali data yang disimpan dalam basis data. & Kriteria: ketepatan dan penguasaan. Bentuk: keaktifan diskusi kelompok; ketepatan jawaban tugas. & Ketepatan perintah SQL-DQL untuk menampilkan kembali data sesuai studi kasus. & 2.5 \\
13 & SCPMK1001-01501 & Perintah SELECT untuk multi table: perintah JOIN dan jenisnya (INNER JOIN, OUTER JOIN, CROSS JOIN). & Bentuk: ceramah, diskusi, latihan analisis. Metode: Case Based Learning (CBL). Penugasan: menuliskan perintah SELECT multi table untuk menampilkan data dari beberapa tabel. & 1 x 2 x 50' tatap muka dan tugas terstruktur; 1 x 2 x 50' tugas mandiri. & Mahasiswa menjelaskan perintah SELECT untuk multi table beserta jenisnya dan menuliskan perintah SELECT multi tabel sesuai kebutuhan. & Kriteria: ketepatan dan penguasaan. Bentuk: keaktifan diskusi kelompok; ketepatan jawaban tugas. & Ketepatan perintah SELECT untuk menampilkan data pada multi table sesuai studi kasus. & 2.5 \\
14 & SCPMK1001-01501 & Kuis 2: materi minggu 10--13. & Ujian teori pilihan ganda dan/atau esai, sifat closed book. & Sesuai jadwal asesmen. & Mahasiswa mengerjakan soal kuis materi minggu 10--13 secara mandiri dan jujur. & Kriteria: ketepatan jawaban. Bentuk: ujian. & Ketepatan jawaban dengan kunci jawaban. & 10 \\
15 & SCPMK1001-01502 & Studi kasus perancangan basis data: identifikasi kebutuhan data (entitas, atribut, relationship, kardinalitas, partisipan) dan perancangan ERD sesuai hasil analisis. & Bentuk: ceramah, diskusi, latihan analisis. Metode: Case Based Learning (CBL). Penugasan: analisis kebutuhan dan perancangan ERD studi kasus. & 1 x 2 x 50' tatap muka dan tugas terstruktur; 1 x 2 x 50' tugas mandiri. & Mahasiswa mengidentifikasi kebutuhan data dan merancang ERD berdasarkan studi kasus kebutuhan organisasi atau bisnis. & Kriteria: ketepatan dan penguasaan. Bentuk: keaktifan diskusi kelompok; ketepatan jawaban tugas studi kasus. & Ketepatan identifikasi kebutuhan data; ketepatan perancangan basis data. & 5 \\
16 & SCPMK1001-01502 & Studi kasus implementasi basis data: pemetaan ERD ke model relasional (3NF), implementasi ke MySQL dengan DDL, manipulasi data dengan DML, dan query SELECT termasuk multi table. & Bentuk: ceramah, diskusi, latihan analisis. Metode: Case Based Learning (CBL). Penugasan: Tugas laporan perancangan dan implementasi basis data. & 1 x 2 x 50' tatap muka dan tugas terstruktur; 1 x 2 x 50' tugas mandiri. & Mahasiswa memetakan rancangan ke model relasional, mengimplementasikan ke DBMS MySQL, memanipulasi data, dan menampilkan data dari beberapa tabel. & Kriteria: ketepatan dan penguasaan. Bentuk: keaktifan diskusi kelompok; ketepatan laporan studi kasus. & Ketepatan query DDL, DML, dan SELECT untuk implementasi dan penampilan data hasil rancangan. & 5 \\
17 & SCPMK1001-01502 & UAS: SQL berkaitan dengan studi kasus, materi minggu 1--16. & Ujian teori pilihan ganda dan/atau esai, sifat closed book. & Sesuai jadwal asesmen. & Mahasiswa mengerjakan soal UAS materi minggu 1--16 secara mandiri dan jujur. & Kriteria: ketepatan jawaban. Bentuk: ujian. & Ketepatan jawaban dengan kunci jawaban. & 20 \\
\end{longtblr}
